\chapter{Lattices and sphere packings}\label{geometry:ch:lattices}

\begin{center}\itshape
\dots the packing will be the tightest possible, so that in no other arrangement could more pellets be stuffed into the same container.
\upshape(J.~Kepler, \emph{The Six-Cornered Snowflake}, 1611, as quoted by T.~C.~Hales)
\end{center}

Sir Walter Raleigh asked his mathematical assistant Thomas Harriot for a formula giving the number of cannonballs in a stacked pile, and
Harriot got interested in the deeper question behind it: what is the best way to stack balls? In a little New Year's essay of 1611, Johannes Kepler, in \emph{On the six-cornered snowflake}, asserted that the greengrocer's stacking of oranges is the tightest possible.
It fills $\frac\pi{\sqrt{18}}\approx74\%$ of space. Gauss proved in 1831 that no \emph{lattice} arrangement does better; Hilbert put the general question
in his eighteenth problem in 1900; and the full proof, by Thomas Hales, needed a huge computer verification and was published in 2005. In 2016 Maryna
Viazovska solved the same problem in dimension $8$ in a short and stunning paper, for which she received the Fields Medal in 2022. The hero of her paper
is the lattice $E_8$, which we construct at the end of this chapter.

On the way we meet the basic tools of the \emph{geometry of numbers}, invented by Hermann Minkowski in the 1890s: a theorem on convex bodies that proves
Fermat's two-squares theorem in four lines, and the reduction of plane lattices, which goes back to Lagrange and Gauss and shows why bees and
greengrocers are right in the plane.

\begin{figure}[ht]
\centering
\fig{geometry-ch09-packings}
\caption{Disks of diameter $1$ centred at the points of the square lattice and of the hexagonal lattice. In red, a fundamental parallelogram: each
contains exactly one disk's worth of area, $\frac\pi4$, and has area $1$ for the square lattice and $\frac{\sqrt3}2$ for the hexagonal one.}
\label{geometry:fig:packings}
\end{figure}

\begin{think}
In Figure~\ref{geometry:fig:packings}, compute the proportion of the plane covered by the disks in each case, using one red parallelogram. Can you imagine a lattice
that does better than the hexagonal one?
\end{think}

\section{Lattices}

\begin{definition}
A \emph{lattice} in $\mathbb R^n$ is a subgroup $\Lambda=\mathbb Zb_1+\dots+\mathbb Zb_n$ generated by a basis $b_1,\dots,b_n$ of $\mathbb R^n$. With $B$ the
matrix whose columns are the $b_i$, $\Lambda=B\mathbb Z^n$. The \emph{Gram matrix} is $G=B^{\mathsf T}B=(b_i\cdot b_j)$, and the \emph{covolume} is
$\det\Lambda=|\det B|=\sqrt{\det G}$.
\end{definition}

\begin{proposition}\label{geometry:prop:latticebases}
Two bases $B$ and $B'$ generate the same lattice if and only if $B'=BU$ for an integer matrix $U$ with $\det U=\pm1$. In particular $\det\Lambda$ does not
depend on the basis.
\end{proposition}

\begin{proof}
If $B\mathbb Z^n=B'\mathbb Z^n$, each column of $B'$ is an integer combination of the columns of $B$, so $B'=BU$ with $U$ an integer matrix; likewise $B=B'V$. Then
$B=BUV$, so $UV=I$ since $B$ is invertible, and $\det U\det V=1$ with both determinants integers: $\det U=\pm1$. Conversely, if $B'=BU$ with such a $U$, then
$U^{-1}$ is an integer matrix (its entries are cofactors divided by $\pm1$), so $U\mathbb Z^n=\mathbb Z^n$ and $B'\mathbb Z^n=B\mathbb Z^n$. Finally
$|\det B'|=|\det B|\,|\det U|=|\det B|$.
\end{proof}

The \emph{fundamental parallelepiped} of a basis is $\Pi=\{t_1b_1+\dots+t_nb_n:0\le t_i<1\}$, of volume $\det\Lambda$. Every $x\in\mathbb R^n$ is uniquely
$\lambda+\pi$ with $\lambda\in\Lambda$ and $\pi\in\Pi$: write $x=\sum t_ib_i$ and split each $t_i$ into its integer and fractional parts. So the translates
$\lambda+\Pi$ tile space, and $\det\Lambda$ is the volume ``per lattice point''. A lattice is \emph{discrete}: a ball contains only finitely many lattice
points, because the coordinates $t=B^{-1}x$ of its points are bounded. So there is a shortest nonzero vector, whose length we write $\lambda_1(\Lambda)$.

\section{Minkowski's theorem}

\begin{lemma}[Blichfeldt]\label{geometry:lem:blichfeldt}
Let $S\subset\mathbb R^n$ be a bounded measurable set with $\operatorname{vol}(S)>\det\Lambda$. Then $S$ contains two distinct points $x,y$ with $x-y\in\Lambda$.
\end{lemma}

\begin{proof}
Cut $S$ along the tiling: $S$ is the disjoint union of the finitely many nonempty pieces $S_\lambda=S\cap(\lambda+\Pi)$. Translate each piece back into $\Pi$:
the sets $S_\lambda-\lambda$ lie in $\Pi$, and their volumes add up to $\operatorname{vol}(S)>\operatorname{vol}(\Pi)$. So two of them overlap: there are
$\lambda\ne\mu$ and a point $z=x-\lambda=y-\mu$ with $x\in S_\lambda$, $y\in S_\mu$. Then $x\ne y$ and $x-y=\lambda-\mu\in\Lambda$.
\end{proof}

\begin{theorem}[Minkowski]\label{geometry:thm:minkowski}
Let $C\subset\mathbb R^n$ be convex, bounded and symmetric ($x\in C\Rightarrow-x\in C$), with $\operatorname{vol}(C)>2^n\det\Lambda$. Then $C$ contains a nonzero
point of $\Lambda$.
\end{theorem}

\begin{proof}
The set $\frac12C$ has volume $2^{-n}\operatorname{vol}(C)>\det\Lambda$, so by Blichfeldt's lemma it contains $x\ne y$ with $x-y\in\Lambda$. Now $2x$ and $-2y$ are in
$C$ (symmetry), so their midpoint $x-y$ is in $C$ (convexity). It is a nonzero point of $\Lambda$.
\end{proof}

The constant $2^n$ cannot be improved: the open cube $(-1,1)^n$ has volume $2^n$ and contains no nonzero point of $\mathbb Z^n$. Here is the classic application.

\begin{theorem}[Fermat]\label{geometry:thm:twosquares}
Every prime $p\equiv1\pmod4$ is a sum of two squares.
\end{theorem}

\begin{proof}
Since $p\equiv1\pmod4$, $-1$ is a square mod $p$ (Volume~1): choose $m$ with $m^2\equiv-1\pmod p$. Let $\Lambda=\{(a,b)\in\mathbb Z^2:b\equiv ma\pmod p\}$, the
lattice with basis $(1,m)$, $(0,p)$ and covolume $p$. The open disk $a^2+b^2<2p$ has area $2\pi p>4p$, so by Minkowski's theorem it contains a nonzero
$(a,b)\in\Lambda$. Then $a^2+b^2\equiv a^2+m^2a^2\equiv0\pmod p$ and $0<a^2+b^2<2p$, so $a^2+b^2=p$.
\end{proof}

\begin{figure}[ht]
\centering
\fig{geometry-ch09-twosquares}
\caption{The proof of Theorem~\ref{geometry:thm:twosquares} for $p=13$, $m=5$: the lattice $b\equiv5a\pmod{13}$ and the disk $a^2+b^2<26$. The four red points are
$\pm(2,-3)$ and $\pm(3,2)$, and $13=2^2+3^2$.}
\label{geometry:fig:twosquares}
\end{figure}

Fermat stated the theorem in a letter to Mersenne of 1640, and Euler gave the first published proof a century later; the lattice proof shows how much a
picture can do. Minkowski's theorem is the foundation of the arithmetic of number fields in Volume~12, Chapter~3. Exercise~\ref{geometry:ex:foursquares} proves
Lagrange's four-squares theorem the same way.

\section{Plane lattices and Gauss reduction}

\begin{definition}
A basis $(b_1,b_2)$ of a plane lattice is \emph{reduced} if $|b_1|\le|b_2|$ and $|b_1\cdot b_2|\le\frac12|b_1|^2$.
\end{definition}

\begin{proposition}\label{geometry:prop:reducedexists}
Every lattice in $\mathbb R^2$ has a reduced basis.
\end{proposition}

\begin{proof}
Let $b_1$ be a shortest nonzero vector of $\Lambda$. In a basis $(e_1,e_2)$, $b_1=pe_1+qe_2$ with $\gcd(p,q)=1$: otherwise $b_1/\gcd(p,q)$ would be a
shorter nonzero vector of $\Lambda$. By B\'ezout there are integers $r,s$ with $ps-qr=1$, and then $(b_1,c)$ with $c=re_1+se_2$ is a basis, by
Proposition~\ref{geometry:prop:latticebases}. Let $k$ be an integer nearest to $\frac{c\cdot b_1}{b_1\cdot b_1}$ and $b_2=c-kb_1$. Then $(b_1,b_2)$ is still a basis,
$|b_1\cdot b_2|=|c\cdot b_1-k\,b_1\cdot b_1|\le\frac12|b_1|^2$, and $|b_2|\ge|b_1|$ because $b_1$ is a shortest vector.
\end{proof}

In practice one finds a reduced basis with Gauss's algorithm (Figure~\ref{geometry:fig:reduction}): subtract from the longer vector the nearest integer multiple of the
shorter one, swap if needed, and repeat until nothing changes. It is the Euclidean algorithm in two dimensions (Exercise~\ref{geometry:ex:gaussreduce}).

\begin{figure}[ht]
\centering
\fig{geometry-ch09-reduction}
\caption{Two bases of the same lattice: a long thin one (red) and the reduced one found by Gauss's algorithm (green). Both parallelograms have the same
area.}
\label{geometry:fig:reduction}
\end{figure}

\begin{theorem}\label{geometry:thm:hermite2}
Let $(b_1,b_2)$ be a reduced basis of a plane lattice $\Lambda$. Then $b_1$ is a shortest nonzero vector, and
\[
\lambda_1(\Lambda)^2\le\frac2{\sqrt3}\det\Lambda,
\]
with equality if and only if $\Lambda$ is hexagonal (generated by two vectors of equal length at an angle of $60^\circ$).
\end{theorem}

\begin{proof}
Let $x=ub_1+vb_2$ with $u,v$ integers, not both zero. If $v=0$, $|x|=|u||b_1|\ge|b_1|$. If $v\ne0$, using $b_1\cdot b_2\ge-\frac12|b_1|^2$ when $uv>0$ and
$b_1\cdot b_2\le\frac12|b_1|^2$ when $uv<0$, together with $|b_2|\ge|b_1|$,
\[
|x|^2=u^2|b_1|^2+2uv\,b_1\cdot b_2+v^2|b_2|^2\ge|b_1|^2\bigl(u^2-|uv|+v^2\bigr).
\]
Now $u^2-|uv|+v^2=\bigl(|u|-\frac{|v|}2\bigr)^2+\frac34v^2\ge\frac34>0$ is an integer, so it is at least $1$, and $|x|\ge|b_1|$. So $|b_1|=\lambda_1$.

For the determinant, $(\det\Lambda)^2=|b_1|^2|b_2|^2-(b_1\cdot b_2)^2\ge|b_1|^4-\frac14|b_1|^4=\frac34\lambda_1^4$, which is the inequality. Equality forces
$|b_2|=|b_1|$ and $|b_1\cdot b_2|=\frac12|b_1|^2$, an angle of $60^\circ$ or $120^\circ$, and both give the hexagonal lattice. Conversely the hexagonal lattice
with $|b_1|=|b_2|=1$ has $\lambda_1=1$ and $\det\Lambda=\frac{\sqrt3}2$, so equality holds.
\end{proof}

In the language of quadratic forms, which is how Lagrange (1773) and Gauss found it: a positive definite form $ax^2+bxy+cy^2$ takes, at some
nonzero integer point $(x,y)$, a value at most $\sqrt{(4ac-b^2)/3}$. For lattices, it gives the optimal packing of equal disks centred at lattice points.

\section{Packings}

If disks (or balls) of radius $r$ are centred at the points of $\Lambda$, they do not overlap as long as $2r\le\lambda_1$. The best choice is
$r=\frac12\lambda_1$, and since there is one ball per fundamental domain, the proportion of space they cover, the \emph{density} of the lattice packing, is
\[
\Delta(\Lambda)=\frac{\operatorname{vol}B\bigl(\frac12\lambda_1\bigr)}{\det\Lambda}=\frac{V_n\lambda_1^n}{2^n\det\Lambda},
\]
where $V_n$ is the volume of the unit ball: $V_2=\pi$, $V_3=\frac43\pi$, $V_4=\frac{\pi^2}2$, $V_8=\frac{\pi^4}{24}$.

\begin{corollary}\label{geometry:cor:hexdensity}
No lattice packing of equal disks in the plane has density larger than $\frac\pi{\sqrt{12}}\approx0.9069$, and only the hexagonal packing reaches it.
\end{corollary}

\begin{proof}
By Theorem~\ref{geometry:thm:hermite2}, $\Delta=\frac{\pi\lambda_1^2}{4\det\Lambda}\le\frac\pi4\cdot\frac2{\sqrt3}=\frac\pi{\sqrt{12}}$, with equality exactly for the hexagonal lattice.
\end{proof}

\begin{remark}
The hexagonal packing is also the densest among \emph{all} packings of equal disks, lattice or not: Axel Thue gave an argument in 1890, and later authors made
it fully rigorous (see Hales's survey in the sources). In dimension $3$, Gauss's 1831 theorem says that the best lattice packing has density
$\frac\pi{\sqrt{18}}$ (Exercise~\ref{geometry:ex:fcc}); that no packing at all does better is Kepler's conjecture, proved by Hales.
\end{remark}

\section{Root lattices and \texorpdfstring{$E_8$}{E8}}

For $n\ge3$, let
\[
D_n=\Bigl\{x\in\mathbb Z^n:x_1+\dots+x_n\ \text{is even}\Bigr\}.
\]
It is a subgroup of index $2$ in $\mathbb Z^n$, hence a lattice with $\det D_n=2$ (Exercise~\ref{geometry:ex:index}). Its shortest vectors are the $\pm e_i\pm e_j$ ($i<j$), of
length $\sqrt2$, and there are $4\binom n2=2n(n-1)$ of them. So
\[
\Delta(D_n)=\frac{V_n(\sqrt2)^n}{2^n\cdot2}=\frac{V_n}{2^{n/2+1}}.
\]
For $n=3$ this is $\frac{4\pi/3}{2^{5/2}}=\frac\pi{\sqrt{18}}\approx0.7405$: $D_3$ is the greengrocer's stacking, the \emph{face-centred cubic} lattice, and each
ball touches $12$ others. For $n=4$ it is $\frac{\pi^2}{16}\approx0.617$, with $24$ neighbours.

In high dimension $D_n$ becomes a poor packing, because it has large holes: the point $h=(\frac12,\dots,\frac12)$ is at distance $\frac{\sqrt n}2$ from every
point of $D_n$ that is nearest to it (Exercise~\ref{geometry:ex:deephole}). For $n=8$ this distance is $\sqrt2$, the minimal distance of $D_8$ itself. So one can
slide a second copy of $D_8$ into the holes without creating shorter distances, and double the density.

\begin{theorem}\label{geometry:thm:e8}
Let $h=(\frac12,\dots,\frac12)\in\mathbb R^8$ and $E_8=D_8\cup(D_8+h)$. Then
\begin{enumerate}
\item $E_8$ is a lattice with $\det E_8=1$;
\item $x\cdot y\in\mathbb Z$ and $x\cdot x$ is even, for all $x,y\in E_8$;
\item the nonzero vectors of minimal length have $x\cdot x=2$, and there are $240$ of them.
\end{enumerate}
\end{theorem}

\begin{proof}
(1) Since $2h=(1,\dots,1)$ has even coordinate sum, $2h\in D_8$. Then $E_8$ is closed under subtraction: the difference of two elements of $D_8+h$ is in $D_8$,
and the difference of an element of $D_8+h$ and one of $D_8$ is in $D_8+h$ or $D_8-h=D_8+h-2h=D_8+h$. So $E_8$ is a subgroup of $\mathbb R^8$ containing
$D_8$ with index $2$ ($h\notin D_8$). It is generated by a basis of $D_8$ and $h$, which spans $\mathbb R^8$; and a subgroup $\Lambda'\supset\Lambda$ with
finite index $k$ is again a lattice, with $\det\Lambda'=\det\Lambda/k$ (Exercise~\ref{geometry:ex:index}). So $\det E_8=\frac22=1$.

(2) For $x\in D_8$, $x\cdot x=\sum x_i^2\equiv\sum x_i\equiv0\pmod2$. For $x=n+h$ with $n\in\mathbb Z^8$ and $\sum n_i$ even,
\[
x\cdot x=\sum_i\Bigl(n_i+\frac12\Bigr)^2=\sum_in_i(n_i+1)+2,
\]
which is even since each $n_i(n_i+1)$ is. Then $2\,x\cdot y=(x+y)\cdot(x+y)-x\cdot x-y\cdot y$ is even, so $x\cdot y\in\mathbb Z$.

(3) A nonzero vector of $D_8$ has $\sum x_i^2\ge2$, with equality for the $4\binom82=112$ vectors $\pm e_i\pm e_j$. A vector of $D_8+h$ has all its coordinates
in $\frac12+\mathbb Z$, so $x\cdot x\ge8\cdot\frac14=2$, with equality exactly when all coordinates are $\pm\frac12$; the condition $\sum n_i$ even (with
$n_i\in\{0,-1\}$) says that the number of minus signs is even, which leaves $2^7=128$ vectors. In total $112+128=240$.
\end{proof}

So $E_8$ is an \emph{even unimodular} lattice, and its packing density is
\[
\Delta(E_8)=\frac{V_8(\sqrt2)^8}{2^8}=\frac{\pi^4}{24}\cdot\frac{16}{256}=\frac{\pi^4}{384}\approx0.2537,
\]
twice that of $D_8$. Each ball touches $240$ others.

\begin{figure}[ht]
\centering
\fig{geometry-ch09-e8}
\caption{The $240$ shortest vectors of $E_8$, projected onto a well-chosen plane (the \emph{Coxeter plane}): they fall on eight circles of $30$ points each.}
\label{geometry:fig:e8}
\end{figure}

\begin{remark}
In 2016 Viazovska proved that no packing of equal balls in $\mathbb R^8$, lattice or not, has density larger than $\frac{\pi^4}{384}$. Her proof constructs a
``magic'' function out of modular forms, which certifies the bound through the linear programming method of Cohn and Elkies. A week later, with Cohn,
Kumar, Miller and Radchenko, she did the same in dimension $24$ for the Leech lattice. Even unimodular lattices like $E_8$ exist only in dimensions divisible
by $8$; Serre's \emph{A Course in Arithmetic} proves this with modular forms. The lattice $E_8$ also appears as the root system of the largest exceptional Lie
group, and Figure~\ref{geometry:fig:e8} is the picture usually drawn of it.
\end{remark}

\section{Exercises}

\begin{exercise}[$\star$]\label{geometry:ex:basisZ2}
Do $(1,2)$ and $(3,4)$ form a basis of $\mathbb Z^2$? Do $(2,3)$ and $(3,5)$? Describe the lattice generated in each case.
\end{exercise}

\begin{exercise}[$\star$]\label{geometry:ex:index}
Let $\Lambda$ be a lattice and $\Lambda'$ a subgroup of $\mathbb R^n$ with $\Lambda\subset\Lambda'$ and $[\Lambda':\Lambda]=k$ finite. Show that $\Lambda'$ is a
lattice with $\det\Lambda'=\det\Lambda/k$. Deduce $\det D_n=2$.
\end{exercise}

\begin{exercise}[$\star$]\label{geometry:ex:gaussreduce}
Reduce the basis $(13,8)$, $(21,13)$ of $\mathbb Z^2$ by Gauss's algorithm. Which famous numbers appear, and what does the algorithm do to them?
\end{exercise}

\begin{exercise}[$\star\star$]\label{geometry:ex:gaussterminates}
Show that Gauss's algorithm terminates and returns a reduced basis.
\end{exercise}

\begin{exercise}[$\star\star$]\label{geometry:ex:foursquares}
(Lagrange) Let $p$ be an odd prime. Show that there are integers $a,b$ with $a^2+b^2+1\equiv0\pmod p$, and that the lattice
$\{x\in\mathbb Z^4:x_3\equiv ax_1+bx_2,\ x_4\equiv bx_1-ax_2\pmod p\}$ contains a nonzero point with $x_1^2+x_2^2+x_3^2+x_4^2<2p$. Deduce that $p$ is a sum
of four squares. (The volume of the ball of radius $r$ in $\mathbb R^4$ is $\frac{\pi^2}2r^4$.)
\end{exercise}

\begin{exercise}[$\star$]\label{geometry:ex:kissing2}
Show that a disk can touch at most $6$ disjoint disks of the same radius.
\end{exercise}

\begin{exercise}[$\star\star$]\label{geometry:ex:fcc}
Show that $D_3$ is generated by $(1,1,0)$, $(1,0,1)$, $(0,1,1)$, that each point has $12$ nearest neighbours, and that the layers $x_3=0$ and $x_3=1$ of $D_3$
are square grids shifted against each other: the greengrocer's pile, seen from the side of a square-based pyramid.
\end{exercise}

\begin{exercise}[$\star\star$]\label{geometry:ex:deephole}
Show that the distance from $h=(\frac12,\dots,\frac12)$ to $D_n$ is $\frac{\sqrt n}2$, and that $D_n\cup(D_n+h)$ is a lattice when $n$ is even, with
minimal norm $\min(2,\frac n4)$. For which $n$ is its density larger than that of $D_n$?
\end{exercise}

\begin{exercise}[$\star\star$]\label{geometry:ex:minkowskibound}
Using Minkowski's theorem with a disk, show that every plane lattice has $\lambda_1^2\le\frac4\pi\det\Lambda$. Compare with Theorem~\ref{geometry:thm:hermite2}.
\end{exercise}

\begin{exercise}[$\star\star$]\label{geometry:ex:circlecount}
Let $N(r)$ be the number of points of $\mathbb Z^2$ in the disk of radius $r$. Show that $|N(r)-\pi r^2|\le 2\sqrt2\,\pi r+2\pi$ for $r\ge1$, so that $N(r)/r^2\to\pi$.
\end{exercise}

\begin{exercise}[$\star$]\label{geometry:ex:e8roots}
Show that the $240$ minimal vectors of $E_8$ span $E_8$, and that the inner product of two distinct minimal vectors $x\ne\pm y$ lies in $\{-1,0,1\}$.
\end{exercise}

\begin{exercise}[$\star\star\star$]\label{geometry:ex:hermite}
(Hermite) Show by induction on $n$ that every lattice in $\mathbb R^n$ has $\lambda_1^2\le\bigl(\frac43\bigr)^{(n-1)/2}(\det\Lambda)^{2/n}$.
\end{exercise}

\section{Hints}

\begin{hint}{geometry:ex:basisZ2}
Compute determinants.
\end{hint}
\begin{hint}{geometry:ex:index}
$k\Lambda'\subset\Lambda$ by Lagrange's theorem, so $\Lambda'$ is discrete. Count lattice points in a large cube.
\end{hint}
\begin{hint}{geometry:ex:gaussreduce}
Each step subtracts the shorter vector from the longer one.
\end{hint}
\begin{hint}{geometry:ex:gaussterminates}
Each step with $k\ne0$ shortens the longer vector; there are only finitely many lattice vectors of bounded length.
\end{hint}
\begin{hint}{geometry:ex:foursquares}
For the congruence, count the values of $a^2$ and of $-1-b^2$ mod $p$. Then mimic Theorem~\ref{geometry:thm:twosquares}; the lattice has covolume $p^2$.
\end{hint}
\begin{hint}{geometry:ex:kissing2}
The centres of two touching neighbours are seen from the centre at an angle of at least $60^\circ$.
\end{hint}
\begin{hint}{geometry:ex:fcc}
The sum of the three generators is $(2,2,2)$; express $e_i\pm e_j$ in terms of them.
\end{hint}
\begin{hint}{geometry:ex:deephole}
The nearest points to $h$ in $\mathbb Z^n$ are the $0/1$ vectors.
\end{hint}
\begin{hint}{geometry:ex:minkowskibound}
A disk of radius $\lambda_1$ contains no nonzero lattice point in its interior.
\end{hint}
\begin{hint}{geometry:ex:circlecount}
Attach to each lattice point the unit square of which it is the lower left corner.
\end{hint}
\begin{hint}{geometry:ex:e8roots}
Use Cauchy--Schwarz and part (2) of Theorem~\ref{geometry:thm:e8}.
\end{hint}
\begin{hint}{geometry:ex:hermite}
Take $b_1$ shortest, project $\Lambda$ orthogonally to $b_1$, and apply induction to the projected lattice, of covolume $\det\Lambda/|b_1|$.
\end{hint}

\section{Solutions}

\begin{solution}{geometry:ex:basisZ2}
\idea A basis of $\mathbb Z^2$ is an integer matrix of determinant $\pm1$.
\writeup $\det\begin{pmatrix}1&3\\2&4\end{pmatrix}=-2$, so $(1,2)$, $(3,4)$ generate a sublattice of index $2$: it has covolume $2$ and is
$\{(a,b):b\ \text{even}\}$, since both generators have even second coordinate and $(1,2)$, $(3,4)-3(1,2)=(0,-2)$ generate it. On the other hand
$\det\begin{pmatrix}2&3\\3&5\end{pmatrix}=1$, so $(2,3)$, $(3,5)$ is a basis of $\mathbb Z^2$.
\end{solution}

\begin{solution}{geometry:ex:index}
\idea A finite-index overgroup is squeezed between $\Lambda$ and $\frac1k\Lambda$.
\writeup For $x\in\Lambda'$, the class of $x$ in the group $\Lambda'/\Lambda$ of order $k$ has order dividing $k$, so $kx\in\Lambda$. Hence
$\Lambda\subset\Lambda'\subset\frac1k\Lambda$, and $\Lambda'$ is a subgroup of the lattice $\frac1k\Lambda=B(\frac1k\mathbb Z^n)$. In coordinates
$t=kB^{-1}x$, $\Lambda'$ becomes a subgroup of $\mathbb Z^n$ of finite index; a subgroup of $\mathbb Z^n$ of finite index has a basis of $n$ vectors (by the
structure theorem for finitely generated abelian groups, or by Hermite normal form), so $\Lambda'$ is a lattice. For the covolume, count points in a large
cube $Q_R=[-R,R]^n$: a lattice $\Gamma$ has $\frac{(2R)^n}{\det\Gamma}+O(R^{n-1})$ points in $Q_R$ (tile by fundamental domains). Since $\Lambda'$ is the union of
$k$ cosets of $\Lambda$, each with $\frac{(2R)^n}{\det\Lambda}+O(R^{n-1})$ points in $Q_R$, comparing the leading terms gives $\frac1{\det\Lambda'}=\frac
k{\det\Lambda}$. For $D_n\subset\mathbb Z^n$: the map $x\mapsto\sum x_i\bmod2$ is a surjective homomorphism $\mathbb Z^n\to\mathbb Z/2$ with kernel $D_n$, so the
index is $2$ and $\det D_n=2\det\mathbb Z^n=2$.
\end{solution}

\begin{solution}{geometry:ex:gaussreduce}
\idea Fibonacci numbers, run backwards.
\writeup With $u=(13,8)$ and $v=(21,13)$: $\frac{u\cdot v}{u\cdot u}=\frac{377}{233}\approx1.618$, rounded to $2$, and $v-2u=(-5,-3)$. Then, with
$u=(-5,-3)$ and $v=(13,8)$, the ratio is $\approx-2.618$, rounded to $-3$, and $v+3u=(-2,-1)$. Next $(-5,-3)-3(-2,-1)=(1,0)$, and finally
$(-2,-1)+2(1,0)=(0,-1)$. The reduced basis is $(1,0)$, $(0,-1)$. The vectors are pairs of consecutive Fibonacci numbers $(F_{k+1},F_k)$, and the ratios are close
to the golden ratio $\varphi\approx1.618$ and to $\varphi^2\approx2.618$: consecutive Fibonacci vectors are nearly parallel, which is why the starting basis is
so bad. With rounding to the nearest integer, each step jumps two places down the Fibonacci sequence; with ordinary division it would go down one place
at a time, the slowest possible run of the Euclidean algorithm (Lam\'e's theorem).
\end{solution}

\begin{solution}{geometry:ex:gaussterminates}
\idea A strictly decreasing sequence of lengths of lattice vectors cannot go on forever.
\writeup Run the algorithm as follows: swap $u$ and $v$ if needed so that $|u|\le|v|$; let $\mu=\frac{u\cdot v}{u\cdot u}$; if $|\mu|\le\frac12$, stop;
otherwise replace $v$ by $v'=v-ku$, with $k$ an integer nearest to $\mu$, and start again. Write $v=\mu u+w$ with $w\perp u$. Then
$|v'|^2=(\mu-k)^2|u|^2+|w|^2$, and when $|\mu|>\frac12$ we have $|\mu-k|\le\frac12<|\mu|$, so $|v'|<|v|$. Each round therefore keeps $u$, replaces $v$ by a
strictly shorter lattice vector, and keeps the basis property (Proposition~\ref{geometry:prop:latticebases}). The sum $|u|^2+|v|^2$ strictly decreases, and all the vectors
involved lie in the ball of radius $|v_0|$, the length of the starting longer vector, which contains finitely many lattice points. So the algorithm stops, and
when it does, $|u|\le|v|$ and $|u\cdot v|\le\frac12|u|^2$: the basis is reduced.
\end{solution}

\begin{solution}{geometry:ex:foursquares}
\idea Same as two squares, in four dimensions, with the quaternion-like congruences making the norm divisible by $p$.
\writeup The $\frac{p+1}2$ values of $a^2$ and the $\frac{p+1}2$ values of $-1-b^2$ mod $p$ cannot be disjoint (together $p+1>p$ values), so $a^2\equiv-1-b^2$ for some
$a,b$. The lattice $\Lambda$ in the statement has the basis given by $x_1,x_2$ free and $x_3,x_4$ determined mod $p$, so its index in $\mathbb Z^4$ is $p^2$ and
$\det\Lambda=p^2$. For $x\in\Lambda$,
\[
x_3^2+x_4^2\equiv(ax_1+bx_2)^2+(bx_1-ax_2)^2=(a^2+b^2)(x_1^2+x_2^2)\equiv-(x_1^2+x_2^2)\pmod p,
\]
so $p$ divides $|x|^2$. The open ball $|x|^2<2p$ has volume $\frac{\pi^2}2(2p)^2=2\pi^2p^2>16p^2=2^4\det\Lambda$, so it contains a nonzero $x\in\Lambda$, and
$0<|x|^2<2p$ with $p\mid|x|^2$ gives $|x|^2=p$. (Every integer is then a sum of four squares, by Euler's four-square identity and $2=1+1+0+0$.)
\end{solution}

\begin{solution}{geometry:ex:kissing2}
\idea Neighbours are at least $60^\circ$ apart as seen from the centre.
\writeup Let the disks have radius $1$, the central one centred at $0$. The centres $c_i$ of touching disks satisfy $|c_i|=2$, and disjointness gives
$|c_i-c_j|\ge2$. In the triangle $0,c_i,c_j$ the side $c_ic_j$ is at least as long as the two others, so it is opposite the largest angle: the angle at $0$ is
at least $60^\circ$. Around $0$ there is room for at most $360/60=6$ such angles.
\end{solution}

\begin{solution}{geometry:ex:fcc}
\idea Every vector with even coordinate sum is an integer combination of the three generators.
\writeup Let $g_1=(1,1,0)$, $g_2=(1,0,1)$, $g_3=(0,1,1)$, all in $D_3$. Then $e_1+e_2=g_1$, $e_1-e_2=g_2-g_3$, and similarly for the other pairs, while $2e_1=g_1+g_2-g_3$.
Every $x\in D_3$ is a sum of vectors $\pm e_i\pm e_j$ (pair off the odd coordinates, whose number is even, and write even coordinates with $2e_i=(e_i+e_j)+(e_i-e_j)$),
so the $g_i$ generate $D_3$; their determinant is $-2=\pm\det D_3$, consistent. The nearest neighbours of $0$ are the $12$ vectors $\pm e_i\pm e_j$. The layer $x_3=0$
is $\{(x_1,x_2,0):x_1+x_2\ \text{even}\}$, a square grid of side $\sqrt2$ (rotated by $45^\circ$), and the layer $x_3=1$ is $\{x_1+x_2\ \text{odd}\}$, the same grid
shifted to sit over the centres of its squares: each ball rests in the hollow between four balls below, as in a pyramid of oranges.
\end{solution}

\begin{solution}{geometry:ex:deephole}
\idea Round $h$ coordinatewise.
\writeup For $x\in\mathbb Z^n$, $|x_i-\frac12|\ge\frac12$ for each $i$, so $|x-h|^2\ge\frac n4$, with equality for the $0/1$ vectors, among which those with an even
number of ones lie in $D_n$. So the distance is $\frac{\sqrt n}2$. If $n$ is even, $2h$ has even coordinate sum, so $2h\in D_n$ and $D_n\cup(D_n+h)$ is a group,
as in Theorem~\ref{geometry:thm:e8}: a lattice of covolume $1$. Its minimal norm is $2$ on $D_n$, and $\frac n4$ on $D_n+h$ (by the computation above, attained), so the
minimum is $\min(2,\frac n4)$. Its density divided by that of $D_n$ is $2\cdot\bigl(\min(2,\frac n4)/2\bigr)^{n/2}$: it equals $2$ for $n=8$, is smaller than $1$ for
$n=4,6$ (where the new points are too close), and for $n>8$ equals $2$ as well.
\end{solution}

\begin{solution}{geometry:ex:minkowskibound}
\idea The open disk of radius $\lambda_1$ is a symmetric convex body with no nonzero lattice point.
\writeup By Minkowski's theorem the open disk of radius $\lambda_1$, which contains no nonzero point of $\Lambda$, has area at most $4\det\Lambda$:
$\pi\lambda_1^2\le4\det\Lambda$. So
$\lambda_1^2\le\frac4\pi\det\Lambda\approx1.273\det\Lambda$, weaker than $\frac2{\sqrt3}\approx1.155$ from reduction.
\end{solution}

\begin{solution}{geometry:ex:circlecount}
\idea Unit squares attached to lattice points tile the plane, and their union is squeezed between two disks.
\writeup To each $(a,b)\in\mathbb Z^2$ attach the square $[a,a+1)\times[b,b+1)$, of diameter $\sqrt2$. The squares of the $N(r)$ points in the disk $D_r$ of radius $r$ are
disjoint, have total area $N(r)$, and lie in $D_{r+\sqrt2}$; and they cover $D_{r-\sqrt2}$ (a point of $D_{r-\sqrt2}$ lies in the square of a lattice point at distance
$<\sqrt2$, which is in $D_r$). So $\pi(r-\sqrt2)^2\le N(r)\le\pi(r+\sqrt2)^2$, and $|N(r)-\pi r^2|\le2\sqrt2\pi r+2\pi$. Dividing by $r^2$, $N(r)/r^2\to\pi$. (The true size of the error is a famous open
problem, Gauss's circle problem.)
\end{solution}

\begin{solution}{geometry:ex:e8roots}
\idea Integrality and Cauchy--Schwarz squeeze the inner products.
\writeup For minimal vectors $x\ne\pm y$, Cauchy--Schwarz gives $|x\cdot y|<|x||y|=2$ (strict, since they are not proportional), and $x\cdot y$ is an integer: so
$x\cdot y\in\{-1,0,1\}$. They span $E_8$: the vectors $e_i\pm e_j$ span $D_8$ (Exercise~\ref{geometry:ex:fcc} for the method), and $h$ itself is a minimal vector (all
coordinates $+\frac12$, zero minus signs), so together they generate $D_8+\mathbb Zh=E_8$.
\end{solution}

\begin{solution}{geometry:ex:hermite}
\idea Peel off the shortest vector and recurse in one dimension less.
\writeup For $n=1$ the claim is $\lambda_1^2\le(\det\Lambda)^2$, an equality. Let $n\ge2$, $b_1$ a shortest vector, and $\pi$ the orthogonal projection onto
$b_1^\perp$. We use two standard facts (Cassels, Chapter~I): $b_1$, being primitive, extends to a basis $b_1,\dots,b_n$ of $\Lambda$; and then
$\pi(\Lambda)$ is a lattice in $b_1^\perp\cong\mathbb R^{n-1}$ with basis $\pi(b_2),\dots,\pi(b_n)$, and $\det\Lambda=|b_1|\det\pi(\Lambda)$
(Gram--Schmidt). By induction there is $y=\pi(x)$, $x\in\Lambda$, with
$|y|^2\le(\frac43)^{(n-2)/2}(\det\Lambda/|b_1|)^{2/(n-1)}$. Changing $x$ by a multiple of $b_1$, we may assume that the component of $x$ along $b_1$ has length at most
$\frac12|b_1|$. Since $x\notin\mathbb Rb_1$, $|x|\ge|b_1|$, so $|b_1|^2\le|x|^2\le\frac14|b_1|^2+|y|^2$, that is $\frac34|b_1|^2\le|y|^2$. Combining,
\[
|b_1|^2\le\frac43\Bigl(\frac43\Bigr)^{\frac{n-2}2}\Bigl(\frac{\det\Lambda}{|b_1|}\Bigr)^{\frac2{n-1}},
\]
so $|b_1|^{2+\frac2{n-1}}=|b_1|^{\frac{2n}{n-1}}\le(\frac43)^{\frac n2}(\det\Lambda)^{\frac2{n-1}}$. Raising to the power $\frac{n-1}n$ gives
$|b_1|^2\le(\frac43)^{\frac{n-1}2}(\det\Lambda)^{\frac2n}$.
\end{solution}

\section*{Sources and further reading}

\begin{itemize}
\item T.~C.~Hales, \emph{A computer verification of the Kepler conjecture}, arXiv:math/0305012 (2003) (Harriot, Raleigh and Kepler's 1611 essay, quoted in the
epigraph); T.~C.~Hales, \emph{A proof of the Kepler conjecture}, Annals of Mathematics 162 (2005), 1065--1185; T.~C.~Hales et al., \emph{A formal proof of the
Kepler conjecture}, Forum of Mathematics, Pi 5 (2017).
\item J.~Kepler, \emph{The Six-Cornered Snowflake} (1611), English translation, Clarendon Press, Oxford, 1966.
\item M.~Viazovska, \emph{The sphere packing problem in dimension 8}, Annals of Mathematics 185 (2017); H.~Cohn, A.~Kumar, S.~D.~Miller, D.~Radchenko and
M.~Viazovska, \emph{The sphere packing problem in dimension 24}, Annals of Mathematics 185 (2017); D.~de~Laat and F.~Vallentin, \emph{A breakthrough in sphere
packing: the search for magic functions}, arXiv:1607.02111 (an account for a general audience).
\item J.~H.~Conway and N.~J.~A.~Sloane, \emph{Sphere Packings, Lattices and Groups}, Springer, 3rd ed., 1999 (everything about $D_n$, $E_8$ and the Leech lattice).
\item H.~Minkowski, \emph{Geometrie der Zahlen}, Teubner, 1896 and 1910; J.~W.~S.~Cassels, \emph{An Introduction to the Geometry of Numbers}, Springer, 1959.
\item J.-P.~Serre, \emph{A Course in Arithmetic}, Springer, 1973 (even unimodular lattices and modular forms).
\item The story of Hilbert's eighteenth problem and of the planar case (Thue) is told in Hales's article above and in C.~Zong, \emph{Sphere Packings},
Springer, 1999.
\end{itemize}
